\section*{5. Griffiths 群}

\subsubsection*{5.0.}
设 $S$ 在代数闭域 $k$ 上固有且光滑. 记

\medskip
\noindent\textbf{(5.0.1)} 以 $Z^i(S)$ 表示 $S$ 上余维为 $i$ 的代数圈群;

\noindent\textbf{(5.0.2)} 以 $Z^i_{\mathrm{alg}}(S)$ 表示 $Z^i(S)$ 中由代数等价于零的圈组成的子群; 并对每个素数 $\ell\ne\operatorname{char}(k)$,

\noindent\textbf{(5.0.3)} 以 $Z^i_{\ell\text{-}\mathrm{coh}}(S)$ 表示 $Z^i(S)$ 中由上同调类在 $H^{2i}(S,\mathbb Q_\ell(i))$ 中为零的圈组成的子群.

\textup{[4]} 中猜想此群与 $\ell$ 无关; 至少当 $\operatorname{char}(k)=0$ 时这是已知的. 置
\begin{equation}\tag{5.0.4}
\mathrm{Grif}^i_\ell(S)
 =Z^i_{\ell\text{-}\mathrm{coh}}(S)/Z^i_{\mathrm{alg}}(S).
\end{equation}
若此外 $S$ 为射影概形, 并配备一个射影嵌入, 则置

\medskip
\noindent\textbf{(5.0.5)} 令 $Z^{i,\mathrm{prim}}(S)$ 为 $Z^i(S)$ 中由本原圈 (即上同调类为本原类的圈) 组成的子群, 并令

\noindent\textbf{(5.0.6)}
$\mathrm{Prim}^{2i}_{\mathrm{alg}}(S,\mathbb Q_\ell(i))$ 为 $H^{2i}(S,\mathbb Q_\ell(i))$ 中由本原代数圈的上同调类张成的 $\mathbb Q$-向量子空间.

\subsubsection*{推论 5.1.}
在定理~4.1 的假设下, 有
\begin{equation}\tag{5.1.1}
\dim_{\mathbb Q}\!\left(
  \mathrm{Grif}^m_\ell(X_{\bar\eta})\otimes_{\mathbb Z}\mathbb Q
\right)
\geq
\dim_{\mathbb Q}\!\left(
  \mathrm{Prim}^{2m}_{\mathrm{alg}}(X,\mathbb Q_\ell(m))
\right).
\end{equation}
事实上, 由 (4.1.5), 复合映射
\begin{equation*}
\mathbb Q\otimes_{\mathbb Z}Z^{m,\mathrm{prim}}(X)
 \xrightarrow{\ \mathrm{res.}\ \text{到 }X_{\bar\eta}\ }
\mathbb Q\otimes_{\mathbb Z}Z^m_{\ell\text{-}\mathrm{coh}}
  (X_{\bar\eta})
 \longrightarrow \mathrm{Grif}^m_\ell(X_{\bar\eta})\otimes_{\mathbb Z}\mathbb Q
\end{equation*}
的核包含于 $Z^m_{\ell\text{-}\mathrm{coh}}(X)\otimes_{\mathbb Z}\mathbb Q$ 中.\EditorialFootnote{译校: 法文印本第 361 页复合映射的末项漏写 $\otimes_{\mathbb Z}\mathbb Q$. 前一项已经有理化, 此处须取商映射的有理化, 与 (5.1.1) 及下一句一致.} 因而 $\mathrm{Prim}^{2m}_{\mathrm{alg}}(X,\mathbb Q_\ell(m))$ 是 $\mathrm{Grif}^m_\ell(X_{\bar\eta})\otimes_{\mathbb Z}\mathbb Q$ 的某个子空间的商, 这就证明了断言.

\subsubsection*{注 5.2.}
若 $X$ 为 $\mathbb P^{2m+1}$ 中的二次超曲面, 则有
\begin{equation*}
\dim_{\mathbb Q}\!\left(
  \mathrm{Prim}^{2m}_{\mathrm{alg}}(X,\mathbb Q_\ell(m))
\right)=1.
\end{equation*}
印本此处将扭写作 $n$, 且漏去了最外层的右括号; 中间次数 $2m$ 要求扭为 $m$.\EditorialFootnote{译校: 法文印本第 362 页还将维数的标量域写作 $\mathbb Q_\ell$, 而 (5.0.6) 定义的是代数圈类张成的 $\mathbb Q$-子空间; 因此按该定义应为 $\dim_{\mathbb Q}$. 对二次超曲面, 两族生成子给出的本原代数圈类之 $\mathbb Q$-空间为一维; 将其扩张为 $\mathbb Q_\ell$-张成空间后也为一维, 见 XII~3.3--3.5.}

特别地, 当 $m=2$ 时, 对于 $\mathbb P^5$ 中的二次超曲面 $X$, $X$ 被次数为 $>5$ 的一般超曲面所截得的截面包含一条曲线, 它在 $\mathbb Q_\ell$-同调中等价于零, 但其任何非零倍数都不代数等价于零.\footnote{印本此处插入了字面括注 ``(4.2)''. 英文版直接查看原始扫描件并作等效于 9000 dpi 的放大, 确认了这些字符; 从 $\mathbb P^{2m+1}$ 到 $\mathbb P^5$ 的数值特化是 $m=2$.}

\begin{thebibliography}{9}
\bibitem{GriffithsPeriodsXX}
P. Griffiths, ``On the Periods of Certain Rational Integrals III,''
即将刊于 \emph{Annals of Mathematics}.

\bibitem{GrothendieckBrauerIIIXX}
A. Grothendieck, ``Le Groupe de Brauer III'' (特别是 \S\S 9--10), 载
\emph{Dix Exposés sur la Cohomologie des Schémas}, North-Holland, 1968.

\bibitem{JouanolouCohomologieXX}
J.-P. Jouanolou, \emph{Cohomologie $\ell$-adique}, 学位论文, Paris, 1969.

\bibitem{KleimanCyclesXX}
S. Kleiman, ``Algebraic Cycles and the Weil Conjectures,'' 载
\emph{Dix Exposés sur la Cohomologie des Schémas}, North-Holland, 1968.

\bibitem{ManinCurvesXX}
J. Manin, ``Rational Points of Algebraic Curves over Function Fields,''
\emph{Translations of the American Mathematical Society}.
\end{thebibliography}
